\section*{How is Claude multilingual?}
Claude speaks dozens of languages fluently-from English and French to Chinese and Tagalog. How does this multilingual ability work? Is there a separate "French Claude" and "Chinese Claude" running in parallel, responding to requests in their own language? Or is there some cross-lingual core inside?\\
\includegraphics[max width=\textwidth, alt={}, center]{multilingual-shared-features}

Shared features exist across English, French, and Chinese, indicating a degree of conceptual universality.

Recent research on smaller models has shown hints of shared grammatical mechanisms across languages. We investigate this by asking Claude for the "opposite of small" across different languages, and find that the same core features for the concepts of smallness and oppositeness activate, and trigger a concept of largeness, which gets translated out into the language of the question. We find that the shared circuitry increases with model scale, with Claude 3.5 Haiku sharing more than twice the proportion of its features between languages as compared to a smaller model.

This provides additional evidence for a kind of conceptual universality-a shared abstract space where meanings exist and where thinking can happen before being translated into specific languages. More practically, it suggests Claude can learn something in one language and apply that knowledge when speaking another. Studying how the model shares what it knows across contexts is important to understanding its most advanced reasoning capabilities, which generalize across many domains.
